\begin{derivation}[从\eqnrefs{eq:maxwell_continuum_delta_frequency_normalization}到\eqnrefs{eq:green_time_response_poles_plus_cut}]
开放无损 Maxwell 算符的谱定理要求离散束缚态与连续散射态一起完备。连续态用频率 $\delta$ 归一，因此把\eqnrefs{eq:maxwell_discrete_continuum_completeness} 插入预解式后得到

\[
\begin{aligned}
\mathbf G^R(\mathbf r,\mathbf r';\omega)
={}&c^2\sum_n
\frac{\mathbf f_n(\mathbf r)\otimes\mathbf f_n^*(\mathbf r')}
{\omega_n^2-(\omega+\ii0^+)^2}\\
&+c^2\sum_a\int_0^\infty\dd\Omega\,
\frac{\mathbf f_{\Omega a}(\mathbf r)\otimes\mathbf f_{\Omega a}^*(\mathbf r')}
{\Omega^2-(\omega+\ii0^+)^2}.
\end{aligned}
\]
第一项给离散极点；第二项因为 $\Omega$ 连续，在复频延拓中产生割线。对 $\omega>0$，Sokhotski--Plemelj 公式给
\[
\operatorname{Im}\frac1{\Omega^2-(\omega+\ii0^+)^2}
=\frac{\pi}{2\omega}\delta(\Omega-\omega),
\]
所以连续谱的在壳谱权重为
\[
\operatorname{Im}\mathbf G_{\rm cont}^R(\mathbf r,\mathbf r';\omega)
=\frac{\pi c^2}{2\omega}
\sum_a\mathbf f_{\omega a}(\mathbf r)
\otimes\mathbf f_{\omega a}^*(\mathbf r').
\]
这正是离散模求和恒等式的连续谱版本。

定义边界值
$\mathbf G^{R/A}(\omega)=\mathbf G(\omega\pm\ii0^+)$，则
\[
\operatorname{Disc}\mathbf G(\omega)
=\mathbf G^R(\omega)-\mathbf G^A(\omega)
=2\ii\operatorname{Im}\mathbf G^R(\omega).
\]
对 $G_j(z)=\langle j|\mathbf G(z)|j\rangle$ 的逆 Fourier 积分轮廓向下半平面闭合。孤立极点给留数，真实频轴上的割线给其上下边界值之差，于是得到\eqnrefs{eq:green_time_response_poles_plus_cut}。该分解也说明：极点主导产生近指数衰减，而阈值或 band-edge 割线可以在长时间产生非指数尾。

若只保留一个目标极点，将投影谱写成
\[
G_j^R(\omega)
=\frac{\mathcal R_j}{\omega-\widetilde\omega_c}
+G_{j,\mathrm{cont}}^R(\omega)+G_{j,\rm an}^R(\omega).
\]
对实际开关窗口 $|\widetilde w(\omega)|^2$ 积分后，割线与解析背景的总权重必须远小于极点权重；这就是已定义的 $\epsilon_{\rm bg}\ll1$ 的谱论来源。靠近波导阈值、光子带边或强辐射连续谱时，这一条件失败，应保留极点+连续谱或直接使用 Green 函数库描述。
\end{derivation}
